\documentclass{article}
\usepackage[T1]{fontenc}
\usepackage{lmodern}
\usepackage{graphicx}
\usepackage{amsmath,amssymb}
\usepackage[a4paper,margin=2cm]{geometry}
\usepackage[absolute,overlay]{textpos}
\setlength{\TPHorizModule}{1mm}
\setlength{\TPVertModule}{1mm}
\usepackage[indonesian]{babel}
\usepackage{hyperref}
\setlength{\emergencystretch}{2em}
% unit: Halaman 15

\begin{document}

% === Halaman 15 (python/native) ===

• Selama kalkulasi plainteks menjadi 

cipherteks, status data sekarang disimpan di 

dalam matriks dua dimensi, state. 

• state bertipe byte dan berukuran Nrows  

Ncols 

• Untuk blok data 128-bit, ukuran state 

adalah 4  4, seperti gambar di samping.

Plainteks 128-bit

state

\end{document}
